\documentclass{article}
\usepackage[T1]{fontenc}
\usepackage{lmodern}
\usepackage{graphicx}
\usepackage{amsmath,amssymb}
\usepackage[a4paper,margin=2cm]{geometry}
\usepackage[absolute,overlay]{textpos}
\setlength{\TPHorizModule}{1mm}
\setlength{\TPVertModule}{1mm}
\usepackage[indonesian]{babel}
\usepackage{hyperref}
\setlength{\emergencystretch}{2em}
% unit: Halaman 25

\begin{document}

% === Halaman 25 (python/native) ===

25

Contoh: kriptogram XMZVH 

Tabel 1. Contoh exhaustive key search terhadap cipherteks XMZVH 


Kunci (k) 

ciphering 

‘Pesan’ hasil 

dekripsi 

Kunci (k) 

ciphering 

‘Pesan’ hasil 

dekripsi 

Kunci (k) 

ciphering 

‘Pesan’ hasil 

dekripsi 

0 

25 

24 

23 

22 

21 

20 

19 

18 

XMZVH 

YNAWI 

ZOBXJ 

APCYK 

BQDZL 

CREAM 

DSFBN 

ETGCO 

FUHDP 

17 

16 

15 

14 

13 

12 

11 

10 

9 

GVIEQ 

HWJFR 

IXKGS 

JYLHT 

KZMIU 

LANJV 

MBOKW 

NCPLX 

ODQMY 

8 

7 

6 

5 

4 

3 

2 

1 

PERNZ 

QFSOA 

RGTPB 

SHUQC 

TIVRD 

UJWSE 

VKXTF 

WLYUG 


Plainteks yang potensial adalah CREAM dengan k = 21. 

Kunci ini digunakan untuk mendekripsikan potongan cipherteks lainnya.

\end{document}
